\section{One actual compatibility stage for all later choices}
\label{sec:common-compatibility-v48}

\status{GCF1--GCF18 construct the original common local comparison stage
on the supplied true-column data.  The output is one \texttt{GJStage},
chosen before every later \(W,V,S,N\), with the full original
\texttt{GIs2}: true volumes, the original numerical local scales, and
\texttt{ColumnDataSplitPrimes}.  In particular all five fields of
\texttt{GroupCompatibilityAll} are constructed together on these
stages.  The proof also preserves the actual comparison map through
arbitrary division-lattice changes.  It retains the original arithmetic
host and column input, and does not manufacture an eventual realization
of an arbitrary already supplied \texttt{FamP} by these stages.
The selected design class already supplies its top estimate; its
remaining family-level issue is the simultaneous realization and its
equality with the given \texttt{Extra.cmp}.}

\paragraph{The precise objects being constructed.}
The accepted source is the pin recorded in this component's provenance.
Write \(E/F\) for the original CM extension and \(a\mapsto\bar a\)
for its fixed involution.  Its finite relative primes are unramified,
as required by the original host.  Let \(\mathbb I_E\) be its idele
group, \(U_E\) the subgroup of finite-component units (with arbitrary
infinite components), and
\[
 U^1(\mathbb A_{F,f})=
 \{z\in\mathbb A_{E,f}^{\times}:z\bar z=1\},\qquad
 U^1(\widehat{\mathcal O}_F)=U^1(\mathbb A_{F,f})
                                  \cap\widehat{\mathcal O}_E^{\times}.
 \tag{GCF.1}
\]
These are exactly the original rank-one unitary group and its integral
subgroup under the finite-adelic base-change identification; in rank
one a matrix is identified with its determinant.  For the true-column
datum \(I\), put \(S=S_G\cup S_Q\), the source's \texttt{uSG}.
Its \texttt{LevOk} says that \(S\cap\bar S=\varnothing\), and that
\(S\cup\bar S\) includes every auxiliary-level prime.  Conversely every prime of
\(S_G=(B.\mathrm{primes}(X)).\mathrm{vGs}\) lies above an
auxiliary level.  Thus \(S\cup\bar S\) is exactly the finite
set where the defining quotient imposes its congruences.
The actual finer determinant group is
\[
 D_{T,Q}=\mathbb I_E/K_{D,Q},\qquad
 U_{O,Q}=\operatorname{im}(U_E\longrightarrow D_{T,Q}).
 \tag{GCF.2}
\]
Here \(K_{D,Q}\) is generated by finite unit ideles congruent to one
at \(S\cup\bar S\), and by the original local \(F^{\times}\)
ideles.  This is the same quotient as HF14--HF15, including its
principal-unit congruences; no quotient by all units is taken until a
relative determinant is formed.

% BEGIN GCF1
\paragraph{GCF1. Hilbert--90 lifts in the actual finite adeles.}
\label{claim:GCF1}
Every \(z\in U^1(\mathbb A_{F,f})\) has an idele lift \(a\in
\mathbb I_E\) satisfying
\[
                         z=a_f/\bar a_f.             \tag{GCF.3}
\]
Two such lifts differ, in their finite components, by an idele of
\(F\).  Every such fixed finite idele, with arbitrary infinite
components, belongs to \(K_{D,Q}\).

At a split place, write \(z=(t,t^{-1})\) and take the lift
\((t,1)\).  At a nonsplit place the elementary field argument applies:
\(a=1+z\) works unless \(z=-1\), in which case choose a nonzero
trace-zero element.  These local choices are units almost everywhere.
For completeness, at any unramified nonsplit place where \(z\) is a
unit, reduce to the quadratic residue-field extension.  Some residue
element \(b\) has \(b+z\bar b\ne0\).  Otherwise putting \(b=1\)
would give \(z=-1\), and the asserted identity for all \(b\) would
make the nontrivial residue involution the identity.  Lift this \(b\)
integrally and set \(a=b+z\bar b\).  It is a unit and satisfies
\(a=z\bar a\).  The split construction is also a unit at every
place where \(z\) is integral.  Thus the local lifts and their inverses
are restricted products, and give a finite idele.  Assigning arbitrary
nonzero infinite components gives (GCF.3).  Dividing two solutions
shows directly that their quotient is fixed by conjugation.

To verify the last assertion with the source's \emph{rational} local
generators, let \(b\) be fixed and consider the finite set consisting
of its nonunit components and the congruence places.  At each underlying
place \(v\) of \(F\), density of \(F\) in \(F_v\) gives
\(k_v\in F^{\times}\) with the same valuation as \(b_v\), and with
\(b_v/k_v\equiv1\) whenever a congruence is required.  Dividing by
the finite product of \(\operatorname{loc}_v(k_v)\) leaves a unit
idele with the prescribed congruences.  The latter and the finite
product both belong to the defining subgroup in (GCF.2).
Infinite components cause no condition in that subgroup.  This also
proves independence of the global lift without changing the original
coefficient field or adding a choice of local representatives to the
quotient.
\hfill\(\square\)
% END GCF1

% BEGIN GCF2
\paragraph{GCF2. The exact congruence criterion for the finer quotient.}
\label{claim:GCF2}
For an idele \(b\), the following two conditions are equivalent to
\(b\in K_{D,Q}\):
\[
 \begin{aligned}
  \operatorname{ord}_P(b_P)&=
                 \operatorname{ord}_{\bar P}(b_{\bar P})
                        &&\text{for every }P,\\
  b_P/(\bar b)_P&\equiv1\pmod P &&\text{for }P\in S.
 \end{aligned}                                                   \tag{GCF.4}
\]
The first equality includes a prime fixed by conjugation.  The second
condition uses the actual conjugation between the two local fields.

Each generator of \(K_{D,Q}\) satisfies (GCF.4).  A local
\(F^{\times}\) idele is fixed; a congruence unit has order zero
and both paired components are one modulo their maximal ideals.
The conditions are multiplicative and closed under inverses, proving
necessity for the subgroup generated by these elements.

Conversely, suppose (GCF.4).  At a place below \(P\in S\), the
extension is split because \(P\ne\bar P\).  Choose
\(k_v\in F^{\times}\) sufficiently close to the \(\bar P\)
component of \(b\).  Its quotient with that component is a principal
unit, and the second equation in (GCF.4) makes the quotient at \(P\)
a principal unit as well.  At a place with no congruence requirement,
choose a local \(F\)-uniformizer of the common order in (GCF.4).
It has that same order upstairs because the original extension is
unramified.  A rational element of \(F\) with this local order can
be used, by density.  Only finitely many places require these choices:
the congruence set is finite, and an idele has finitely many nonunit
components.  Dividing \(b\) by their finite local product leaves an
element of the defining congruence-unit subgroup.  This proves
sufficiency.

In particular, forgetting the congruences gives the original identity
\[
 K_{D,Q}U_E=K_DU_E
 =\{b:\operatorname{ord}_P b_P
                  =\operatorname{ord}_{\bar P}b_{\bar P}
                    \text{ for every }P\}.             \tag{GCF.5}
\]
This agrees with HF15 and will preserve its positive-part relative
determinant representative.
\hfill\(\square\)
% END GCF2

% BEGIN GCF3
\paragraph{GCF3. The canonical norm-one determinant map, with all scalar laws.}
\label{claim:GCF3}
Define
\[
 d_Q:U^1(\mathbb A_{F,f})\longrightarrow D_{T,Q},\qquad
                d_Q(z)=[a]\quad(a_f/\bar a_f=z).        \tag{GCF.6}
\]
GCF1 makes this well defined, and products of lifts show it is a
homomorphism.  It is surjective: for any representative \(a\), use
\(z=a_f/\bar a_f\).  Under the rank-one matrix dictionary, (GCF.6)
is the original \texttt{u1Dt}; it is characterized by
\(d_Q(z)=[a]\) for \emph{every} \texttt{IsDetLift}, not merely for a
selected lift.

For a lift \(a\),
\[
 \operatorname{ord}_P z_P=
       \operatorname{ord}_P a_P-
                  \operatorname{ord}_{\bar P}a_{\bar P}.
 \tag{GCF.7}
\]
Thus \(z\) is integral if and only if the orders of \(a\) are
balanced: integrality of a norm-one element includes integrality of
its inverse, since the inverse is its conjugate.  By GCF2,
\[
 \ker d_Q=
 \{z\in U^1(\widehat{\mathcal O}_F):z_P\equiv1\pmod P
                                                    \ (P\in S)\}.
 \tag{GCF.8}
\]
This is open in the actual finite-adelic topology.  Indeed the integral
unit subgroup is open in the restricted product of local unit groups,
and finitely many principal-unit conditions define an open subgroup.
The identification \(E\otimes_F\mathbb A_{F,f}\simeq
\mathbb A_{E,f}\) has this topology: decompose a local tensor product
by the primes above \(v\); its integral subring is the product of the
integral closures, and then take the restricted product.  Consequently
(GCF.8) supplies the source's \texttt{dt\_ker} and \texttt{dt\_open}
on its specified tensor-product model.

The image of the integral subgroup is exactly \(U_{O,Q}\).
One inclusion follows by taking a unit idele lift.  For the reverse
inclusion, any integral \(z\) has a lift with balanced orders by
(GCF.7); (GCF.5) writes its class as the class of a unit idele.
Finally, for rational \(z\in E^{\times}\) of norm one, choose
\(k\in E^{\times}\) with \(z=k/\bar k\).  The explicit solutions
in HF14 give the original rational determinant map
\(e_{1,Q}(z)=[\operatorname{diag}k]\).  The diagonal lift in
(GCF.6) therefore proves
\[
 d_Q(z_f)=e_{1,Q}(z),\qquad
 d_Q\bigl(U^1(\widehat{\mathcal O}_F)\bigr)=U_{O,Q}.    \tag{GCF.9}
\]
These are all six scalar determinant fields: the map, surjectivity,
kernel inclusion, open kernel, integral-unit image, and rational-point
law.  The last law also covers any rank-one unitary matrix whose
determinant is the specified rational element.
\hfill\(\square\)
% END GCF3

% BEGIN GCF4
\paragraph{GCF4. Reading the determinant of an actual decorated point.}
\label{claim:GCF4}
Let \(p\) be an original framed decorated self-dual lattice of rank
\(n\), let \(g\) realize its extension to the prescribed \(S_Q\)
frames, and let \(a\) be a Hilbert--90 lift of \(\det g\).  Then
\[
                         [a]=\operatorname{detPt}(p)
                                       \quad\text{in }D_{T,Q}.
 \tag{GCF.10}
\]
Here every frame is the one in the given decoration.  Choose its
rational local basis matrix \(F_P\), and put
\(d_P=\operatorname{ord}_P\det F_P\).  The local duality argument
in HF15 gives \(d_{\bar P}=-d_P\).  The source's point idele
\(a_p=\operatorname{ptIdl}(p)\) has components
\[
 (a_p)_P=
 \begin{cases}
    \det F_P,&P\in S,\\
    1,&\bar P\in S,\\
    \pi_P^{\max(d_P,0)},&P,\bar P\notin S.
 \end{cases}                                                   \tag{GCF.11}
\]
The set \(S\) contains no conjugate pair, so these cases are
unambiguous.  They give
\(\operatorname{ord}_P(a_p)_P-
\operatorname{ord}_{\bar P}(a_p)_{\bar P}=d_P\).
The actual carried-lattice equality gives
\(\operatorname{ord}_P\det g_P=d_P\), hence \(a/a_p\) has
balanced orders by (GCF.7).

At a selected prime put \(K_P=g_P^{-1}F_P\).  Both matrices carry
the standard local lattice onto the given local lattice; therefore
\(K_P\) and its inverse are integral.  The residue-frame condition
in \texttt{IsPtOf} gives \(K_P\equiv1\pmod P\) at an auxiliary
prime.  At a prescribed rational-frame prime it gives \(K_P=1\).
These assertions can be checked column by column after multiplying
each rational local basis vector by a denominator outside \(P\).
Reduction of the determinant consequently gives
\(\det K_P\equiv1\pmod P\).  Since
\(\det g_P\det K_P=\det F_P\), (GCF.11) implies
\[
 \frac{(a/a_p)_P}{\overline{(a/a_p)}_P}
                 =(\det K_P)^{-1}\equiv1\pmod P.
\]
GCF2 proves \(a/a_p\in K_{D,Q}\), establishing (GCF.10).

Now suppose the original \texttt{Reads} datum realizes the class
by \(y k\), where \(k\) is unitary, preserves the standard lattice,
and is the identity at all the selected underlying places.  Its
determinant is integral and is one at those places, so GCF3 gives
\(d_Q(\det k)=1\).  Therefore (GCF.10) becomes
\[
                \operatorname{detPt}(p)=d_Q(\det y).
 \tag{GCF.12}
\]
For the reference tuple, the twist-prime set is empty by its literal
definition.  The \texttt{Reads.k\_off} clause therefore supplies this
identity condition at the auxiliary and \(S_Q\) places; the type
primes are disjoint from them.  This verifies the original
\texttt{pointDeterminant\_eq\_}\allowbreak\texttt{normOneClass\_of\_level} on the actual
class representatives, including rank zero, where all determinants
are one.
\hfill\(\square\)
% END GCF4

% BEGIN GCF5
\paragraph{GCF5. All column determinant readings and the adelic relative law.}
\label{claim:GCF5}
Take the actual \(C:\mathrm{ColsTiedSym}(I.\mathrm{rec7}(l,Z))\)
and its original \(\mathrm{StdC}\).  On its underlying Hermitian
columns, use \(d_Q\) from GCF3.  The six scalar fields in that claim
and the following two readings give the full
\texttt{SymmetricHermitianColumns} record:
\[
 \begin{aligned}
 D.\mathrm{detCls}(u,0,c)&=d_Q(\det y_{u,0,c}),\\
 D.\mathrm{detCls}(u,i,c)&=d_Q(\det y_{u,i,c})
                                  &&(1\le i\le m).
 \end{aligned}                                                   \tag{GCF.13}
\]
For step zero use the endpoint \texttt{Reads} in \(\mathrm{StdC}\)
and the \texttt{Car0} representative equality.  For
\(1\le i<i+1\le m\), use its middle \texttt{Reads} and the
corresponding chart carrier.  For \(i=m\), use its terminal
\texttt{Reads}, the singleton terminal chart, and \texttt{CarM}.
Each is precisely (GCF.12).  The carriers only reindex the matrix
basis by the supplied rank equality; a simultaneous row and column
reindexing leaves the determinant unchanged.  Thus (GCF.13) holds for
the original representatives and for every rank-one unitary matrix
having their determinant, as demanded by \texttt{det0} and
\texttt{detI}.

There is also the full relative law for \emph{every} rank \(N\) and
every pair of self-dual lattices \(M,M'\).  Suppose an actual unitary
finite adelic matrix \(g\) carries \(M\) onto \(M'\), and a
rank-one unitary element \(z\) has determinant \(\det g\).  Choose
local basis matrices for the two lattices.  In those bases the
carried-lattice condition is an invertible integral matrix, so
\[
 \operatorname{ord}_P\det g_P=d_P(M')-d_P(M)=:\Delta_P.
 \tag{GCF.14}
\]
This uses local bases only.  It does not presume either global lattice
is free.  Completion of a finite lattice and localization give the
same integral span at \(P\); this verifies that the source's
\texttt{Carries} has exactly the interpretation used in (GCF.14).
For a lift \(a\) of \(z\), (GCF.7) gives the same differences
\(\Delta_P\).  The source relative idele is the positive-part
product in HF15,
\[
              a_{M,M'}=\prod_P
                  \operatorname{loc}_P(\pi_P)^{\max(\Delta_P,0)}.
\]
Self-duality gives \(\Delta_{\bar P}=-\Delta_P\), so this product
has the same conjugate-order differences as \(a\).
Their ratio belongs to \(K_{D,Q}U_E\) by (GCF.5).  Passing to the
quotient by \(U_{O,Q}\) proves exactly
\[
       D.\operatorname{relDet}(N,M,M')=[d_Q(z)]
                                  \quad\text{in }D_{T,Q}/U_{O,Q}.
 \tag{GCF.15}
\]
This is \texttt{SymmetricRelativeDeterminant}, with its original
positive-part convention and no rationality restriction on \(g\).
\hfill\(\square\)
% END GCF5

% BEGIN GCF6
\paragraph{GCF6. The full Hermitian extension preserves the given local ties.}
\label{claim:GCF6}
For the original \(I,\mathrm{Ok},\mathrm{Pins},l,Z,C,\mathrm{StdC}\),
suppose true volume compatibility is supplied and fix an actual
\(g:\mathrm{SymmetricGroupCompatibility}(D,C.\mathrm{herm.base})\).
There exists the source's \(F_u:\mathrm{ColsFullS}(C)\) such that
\[
 F_u.\mathrm{g}=g,\qquad
 I.\mathrm{StdDt}(l,Z,F_u.\mathrm{base.dt}).             \tag{GCF.16}
\]
The first equality is equality after the original dependent
identifications; it preserves the supplied local ties themselves.

Extend \(C.\mathrm{herm.base}\) by the eight fields proved in
GCF3 and GCF5.  Take its \texttt{atT}, level-zero compatibility,
adelic compatibility, cofactors, auxiliary-generator identity,
stabilizer maps, residue-field levels, standard lattice and base
compatibility directly from the supplied \(C.\mathrm{herm}\).
Insert the given \(g\), the supplied volume identity, (GCF.15), and
the original positive-volume field \(B.V_{\rm pos}\).
This lists every field of \texttt{FullTiedS}.  Forgetting the new
determinant, volume and local-group fields returns exactly
\(C.\mathrm{herm}\), which is the defining equality in
\texttt{ColsFullS}.  The chosen group field is still \(g\).
Finally the \texttt{StdDt} assertion quantifies over every idele lift
of every rank-one element; it is exactly the universal characterization
of (GCF.6).  This proves the entire original
\texttt{exists\_fullHermitianColumns\_trueColumns} endpoint.
\hfill\(\square\)
% END GCF6

\subsection{The given period root needs no new coefficient field}

% BEGIN GCF7
\paragraph{GCF7. The finite torsion action with its exact scaling factor.}
\label{claim:GCF7}
Let \(\Lambda\) be a complex period lattice and write
\(G=g_2(\Lambda)\), \(T=g_3(\Lambda)\).  Suppose an automorphism
\(\sigma\) of \(\mathbb C\) fixes \(G^3\) and \(T^2\).
There is \(\eta\in\mathbb C\), \(\eta^6=1\), with
\[
 \sigma G=\eta^2G,\qquad \sigma T=\eta^3T.             \tag{GCF.17}
\]
For every torsion class \(t\pmod\Lambda\) there is a class
\(t'\pmod\Lambda\) such that, for every scalar \(c\) preserving
\(\Lambda\) and fixed by \(\sigma\),
\[
 \begin{split}
 ct\in\Lambda&\ \Longleftrightarrow\ ct'\in\Lambda,\\
 \sigma\wp_\Lambda(ct)&=\eta\wp_\Lambda(ct')
                                      \quad(ct\notin\Lambda).
 \end{split}                                                   \tag{GCF.18}
\]

If \(G\ne0\), put \(u=\sigma G/G\), so \(u^3=1\);
if \(G=0\), take \(u=1\).  The equality for \(T^2\) gives
\(\sigma T=T\) or \(-T\).  Accordingly choose
\(\eta=u^2\) or \(-u^2\).  These choices prove (GCF.17), including
the zero-invariant cases.  Choose \(m\) with \(m^2=\eta^{-1}\).
Then \(m\ne0\), and the scaling laws PA1 give
\[
 g_2(m\Lambda)=\sigma G,\qquad
 g_3(m\Lambda)=\sigma T,\qquad m^{12}=1.
\]
Apply coordinatewise \(\sigma\) to the elliptic curve
\(y^2=x^3-Gx/4-T/4\), and uniformize the image by
\(\mathbb C/(m\Lambda)\).  TA1 identifies precisely this
conjugate lattice, and TA5 proves that addition, torsion and CM
multiplication are transported by their algebraic formulas.  A
CM endomorphism with differential \(c\) becomes one with differential
\(\sigma c=c\).  Dividing the image point's parameter by \(m\)
therefore gives \(t'\) with the first assertion in (GCF.18).
The identity
\(\wp_{m\Lambda}(mz)=m^{-2}\wp_\Lambda(z)=\eta\wp_\Lambda(z)\)
gives the second.  This argument applies a possibly discontinuous
\(\sigma\) to algebraic curve identities, not to analytic limits or
complex Taylor series.
\hfill\(\square\)
% END GCF7

% BEGIN GCF8
\paragraph{GCF8. The distinguished two-torsion point fixes the sign of \(g_3\).}
\label{claim:GCF8}
Suppose \(G^3-27T^2=1\), and a scalar \(w\) preserves
\(\Lambda\), satisfies \(w(1-w)=2k\) for \(k\in\mathbb Z\),
and admits a class \(t\) with
\[
       2t\in\Lambda,\qquad wt\notin\Lambda,\qquad
                                      (1-w)t\notin\Lambda.
 \tag{GCF.19}
\]
Every complex field automorphism fixing \(w,G^3,T^2\) fixes
\(T\) itself.

By GCF7 it transports \(t\) to a class \(t'\) with exactly
(GCF.19).  TA11 proves \(t'=t\pmod\Lambda\): on the
two-dimensional \(\mathbb F_2\)-space of two-torsion, \(w\) and
\(1-w\) are complementary idempotents, and the class with both
components nonzero is unique.  Put
\(x=\wp_\Lambda(wt)\) and \(y=\wp_\Lambda((1-w)t)\).
Both are abscissae of nonzero two-torsion points and satisfy
\[
                 4X^3-GX-T=0.                        \tag{GCF.20}
\]
They are distinct: equality of two \(\wp\)-values means equality
up to sign modulo the lattice; since both classes are two-torsion,
this would make their sum \(t\) zero and contradict (GCF.19).
Subtracting their two equations gives
\(G=4(x^2+xy+y^2)\).  Substitution into either equation gives
\[
 \{4(x-y)(2x+y)(x+2y)\}^2=G^3-27T^2=1.               \tag{GCF.21}
\]
The factor in braces is therefore \(1\) or \(-1\), so \(\sigma\)
fixes it.  On the other hand (GCF.18) and uniqueness of \(t\)
multiply each of \(x,y\) by \(\eta\); the same factor is multiplied
by \(\eta^3\).  Thus \(\eta^3=1\).  Equation (GCF.17) proves
\(\sigma T=T\).

The original CM ideal lattices meet the existence clause in (GCF.19).
Here \(w=(1+i\sqrt d)/2\), and the host gives
\(d\equiv7\pmod8\), so \(w(1-w)=(d+1)/4\) is even.
For the unit ideal PA2 gives \(t=\delta/2\) directly.
For an arbitrary nonzero ideal \(\mathfrak a\subset\mathcal O_K\),
\(\mathcal O_K/2\simeq\mathbb F_2\times\mathbb F_2\), with
\(w\) the two complementary idempotent coordinates.  Invertibility
of \(\mathfrak a\) makes \(\mathfrak a/2\mathfrak a\) a
rank-one module over each factor.  Its unique element nonzero in
both factors, divided by two and multiplied by \(\delta\), gives
the required class on \(\delta\iota(\mathfrak a)\).
This verifies the source's half-point input without extending the
coefficient number field.
\hfill\(\square\)
% END GCF8

% BEGIN GCF9
\paragraph{GCF9. Three-torsion controls the cubic ambiguity of \(g_2\).}
\label{claim:GCF9}
Suppose \(G^3-27T^2=1\), \(\beta\Lambda\subset\Lambda\),
\(\beta^2=-d\), where \(d>0\) and \(3\nmid d\).
There is \(v\in\mathbb C\), \(v^3=1\), such that every
automorphism fixing \(\beta,G^3,T^2\) fixes \(G/v\).

First choose a three-torsion class \(t\) not killed by either
\(\beta-1\) or \(\beta+1\).  Multiplication by \(\beta\)
has an integral matrix of trace zero on a basis of the lattice:
its complex eigenvalues are \(\beta,\bar\beta=-\beta\).
On \(\Lambda/3\Lambda\simeq\mathbb F_3^2\), neither
\(\beta-1\) nor \(\beta+1\) is the zero map, since the
trace of \(1\) or \(-1\) is nonzero in \(\mathbb F_3\).
Their two kernels have at most three elements each and cannot cover
the nine-element space.  Choose a class outside their union and
divide its lattice representative by three.  This gives
\[
 3t\in\Lambda,\qquad(\beta-1)t\notin\Lambda,\qquad
                              (\beta+1)t\notin\Lambda.          \tag{GCF.22}
\]
Multiplication by \(\beta\) is invertible on three-torsion, since
\(3\nmid d\).  Its square is \(1\) or \(-1\) there, so it
acts as an involution on the four nonzero abscissae, which identify
opposite classes.  By (GCF.22),
\(x=\wp(t)\) and \(y=\wp(\beta t)\) form a two-element orbit.

Every nonzero three-torsion abscissa \(X\) satisfies
\[
 3X^4-\tfrac32GX^2-3TX-G^2/16=0.                    \tag{GCF.23}
\]
For an explicit derivation, on \(y^2=x^3-Gx/4-T/4\) the tangent
slope is \((3x^2-G/4)/(2y)\).  A nonzero three-torsion point has
\(y\ne0\) and \(2P=-P\), so its doubled abscissa is \(x\).
The duplication formula gives the slope squared equal to \(3x\);
substitution yields (GCF.23).  The four opposite pairs have distinct
abscissae by the degree-two \(\wp\) map, so they are exactly the
four roots of (GCF.23).  Their sum is zero.

For two distinct roots \(x,y\), subtracting their equations yields
\(T=(x+y)(x^2+y^2)-G(x+y)/2\).  Subtracting their equations after
multiplication by the other root also yields
\[
 -3xy\{(x+y)^2-xy\}+\tfrac32Gxy-G^2/16=0.
\]
Eliminating \(T\) with these two identities gives
\[
                   \{G-3(x+y)^2\}^3=G^3-27T^2=1.     \tag{GCF.24}
\]
Set \(v=G-3(x+y)^2\), which is therefore nonzero.

Apply GCF7 to an automorphism fixing the stated data.  It sends the
chosen orbit to another two-element orbit \(\{x',y'\}\) of the
same involution.  The two orbits either agree or are disjoint.
In the latter case they are all four roots of (GCF.23), so
\(x'+y'=-(x+y)\).  In both cases
\((x'+y')^2=(x+y)^2\).  Consequently
\[
 \sigma v=\eta^2v,\qquad
                 \sigma(G/v)=G/v,                   \tag{GCF.25}
\]
as claimed.  Repeated roots or an unspecified pairing are not used:
the opposite three-torsion classes give the four distinct roots and
the original CM scalar gives their involution.
\hfill\(\square\)
% END GCF9

% BEGIN GCF10
\paragraph{GCF10. The prescribed twelfth root has invariants in the given field.}
\label{claim:GCF10}
Take the original \(D\), \(\mathrm{FieldOK}\), and a prescribed
\(\delta\) satisfying
\(\delta^{12}=\Delta_{\iota(\mathrm{ccK})}\).
Let \(M=\operatorname{im}(D.\mathrm{embLL})\subset\mathbb C\),
and \(\Lambda=\delta\iota(\mathrm{ccK})\).
PA1 gives \(\Delta_\Lambda=1\), and the first \texttt{FieldOK}
clause and homothety invariance of \(j\) give
\[
  G^3=j(\iota(\mathrm{ccK}))/1728\in M,\qquad
                    T^2=(G^3-1)/27\in M.             \tag{GCF.26}
\]
The image of the original CM field is in \(M\), so it contains
\(w=(1+i\sqrt d)/2\) and \(\beta=2w-1\).
The host has \(3\nmid d\).  GCF8 and GCF9 show that all
automorphisms of \(\mathbb C\) fixing \(M\) fix \(T\) and
\(G/v\), for a fixed \(v^3=1\).  Hence
\[
                              T\in M,\qquad G/v\in M.
 \tag{GCF.27}
\]
The fixed-field assertion here is the one used in TA5: an algebraic
element outside a characteristic-zero field has a different conjugate,
and an element transcendental over it can be moved by extending a
transcendence basis.  Either embedding extends to an automorphism of
\(\mathbb C\).  Thus an element fixed by all these automorphisms is
in \(M\).

For arbitrary \(M\) this already proves the original
\texttt{exists\_delta\_data}: put \(u=v^{-1}\) and
\(\delta'=u^2\delta\).  Then \((\delta')^{12}=\delta^{12}\),
\[
             g_2(u^2\Lambda)=uG=G/v\in M,\qquad
                         g_3(u^2\Lambda)=T\in M.       \tag{GCF.28}
\]
For the actual \(\mathrm{Base}\), no period is changed: its
coefficient field contains all third roots of unity, since its
prescribed roots field contains \(\mathbb Q(\mu_{3^4})\).
Thus \(v\in M\), and (GCF.27) gives \(G,T\in M\) for the
\emph{given} \(\delta\).  Equivalently one may compare it with
\(\delta'\): their ratio has twelfth power one, so its fourth
power is a third root of unity and its sixth power is \(\pm1\).
The scaling factors for \(g_2,g_3\) already lie in \(M\).

All remaining \texttt{DataIn} coordinates are
\(D.\mathrm{embLL}(D.\mathrm{embE}(\sigma_\tau a))\), and
therefore also lie in \(M\).  This supplies the complete original
\texttt{DataIn} for \(\mathrm{tor0}(D,\mathrm{ccK},\delta)\)
without enlarging \(D.\mathrm{LL}\), changing its embedding, or
renormalizing the originally chosen period root.
\hfill\(\square\)
% END GCF10

% BEGIN GCF11
\paragraph{GCF11. The original \(\Omega\) formula holds for the true global stage.}
\label{claim:GCF11}
Use GCF10 on the original base with its third roots of unity, and
write \(e:D.\mathrm{LL}\simeq M\) for its image-field isomorphism.
For every original infinity label \(\tau\), let \(d_\tau\in M\)
be the image of the source's chosen generator
\(\mathrm{gen}(D.\mathrm{host},D.\mathrm{embE},(D.\mathrm{eTau}(\tau))_1)\).
The second \texttt{FieldOK} clause gives
\[
                 \ell_T(D,M,\tau)=\mathcal O_Md_\tau.
 \tag{GCF.29}
\]
Indeed the source's \(\ell_T\) is generated by the images of
\(a b\), where \(a\in A\) and \(b\) belongs to the original
different ideal, with the stated embedding at \(\tau\).
This is exactly \(\mathrm{ellLat}\).  The generator identity in
the coefficient field carries to (GCF.29) under the isomorphism
\(e\); in each direction, express an element as a finite integral
linear combination and map every coefficient.  The ideal is nonzero,
so \(d_\tau\ne0\).

For \(s:D.\mathrm{LL}\hookrightarrow\mathbb C\), put
\(\rho=s\circ e^{-1}\).  The source's \(\mathrm{conjLat}\)
and the specified \(\mathrm{lamConj}(M,\rho,\Lambda)\) both
have discriminant one and the same \(j\)-invariant.  PA1 and TA1
then show they differ by a homothety \(\zeta\) with
\(\zeta^{12}=1\).  To check the last normalization, equality of
\(j\) and of the discriminants gives equality of the cubes of
\(g_2\) and squares of \(g_3\).  The scalar calculation in GCF7
provides a simultaneous homothety, also when one invariant is zero;
uniqueness of the lattice with its two invariants is TA1.
Its absolute value is one.  Their areas are therefore equal.

The original \texttt{trueGx} is defined using the area of
\(\mathrm{conjLat}\) and the \emph{same} chosen generator.
Its positive-value guard is inactive, since the area is positive
and \(\rho(d_\tau)\ne0\).  Thus it obeys the exact identity
\[
 \mathrm{trueGx}.\Omega(\tau,s)
   =\left\{
       \frac{\pi\,|\rho(d_\tau)|^2}
       {\operatorname{area}(\mathrm{lamConj}(M,\rho,\Lambda))}
     \right\}^{-1}.                                  \tag{GCF.30}
\]
The original \(\mathrm{OmegaCompatible}\) consists precisely of
the image field \(M\), its isomorphism \(e\), its given complex
embedding, \texttt{DataIn}, and (GCF.29)--(GCF.30) for every
\(\tau,s\).  GCF10 supplies \texttt{DataIn}; the preceding argument
supplies every other field.  Consequently \(\mathrm{retie}(D)\)
has this complete property.  If the original later stage already has
\(G.\mathrm{gx}=\mathrm{trueGx}\), the formula holds for that
stage itself.  In the selected application \(\mathrm{ccK}=\mathcal
O_K\), and the ambient input \(I.J.\delta\) is exactly the period
root just used.  This proves the literal
\texttt{omTied\_withTrue}, for all later \(W,V,S,N\).
\hfill\(\square\)
% END GCF11

\subsection{Constructing and preserving the same local comparison}

% BEGIN GCF12
\paragraph{GCF12. All five compatibility fields on the actual recursive record.}
\label{claim:GCF12}
Fix a supplied scale \(\mathrm{sc}\), original level-zero data
\(L_0\), forms \(T_u\), \texttt{StdTied}, \texttt{Joint0},
coefficient-field level-zero characters, and the original companion
identity
\[
              L_{b,u}+\mathfrak a_u L_{b,u}^{\#}=\Lambda_0.
 \tag{GCF.31}
\]
There is one family \(g_{u,w}\) of the actual local models, with
their exact rank identities, for which replacing the \texttt{gat}
field gives an original record \(D^g\), together with
\[
 G_t:\mathrm{GatTies}(D^g),\quad c:\mathrm{LatCmp}(D^g),\quad
 \mathcal G:\mathrm{GroupCompatibilityAll}(D^g,G_t,c).
 \tag{GCF.32}
\]
Its level-zero datum is the given \(L_0\); its forms are the given
\(T_u\), transported to that record.  It also satisfies the full
\texttt{IsRecG}, including its joint equality of
\((D^g,G_t,c,\mathcal G.\mathrm{herm},\mathcal G.\mathrm{forms})\).

Here are the actual choices and all the fields.  LI1--LI8 construct
the local completion identifications \(I_{u,w}\), inverse away
characters \(\psi_{u,w}\), and the embeddings into the same
\(\mathbb C_p\).  SF1--SF7 use (GCF.31) to choose one
\(B\in\mathrm{GL}_k(E)\) with
\[
 B^*JB=1,\qquad
 x\in\Lambda_0\ \Longleftrightarrow\ B^{-1}x\in\mathcal O_E^k,
 \tag{GCF.33}
\]
and make every local frame the specified base change of that \(B\).
Take the original \(\mathrm{packC}\) and \(\mathrm{gatOf}\)
with these choices, with \(C_e=1\), and with the original
\(r_{\mathrm{sc}}\) and \(\mathrm{ep}_{\mathrm{sc}}\).
CR1--CR6 give the exact equality of its local coefficient rank
with \(\dim_{\mathbb C}(D.\mathrm{mem}(u).A.Mu0)\).
Consequently replacing the \texttt{gat} field is legitimate and is
exactly \(\mathrm{recC}(D,\mathrm{sc},1,L_0,I,\psi,h_d)\).

Take \(\mathcal G.\mathrm{herm}\) to be \(L_0\) with the
canonical local ties of this construction, and take the transported
\(T_u\) as \(\mathcal G.\mathrm{forms}\).  VE8--VE9 provide all
six original \texttt{CmpHyp0P} fields on these same objects:
polynomial value existence, evaluation injectivity, point existence,
all presentations, inverse-character identity, and integral lattice
span.  Define \(c\) by those unique value polynomials and the
original coefficient embedding.  The exact equality is
\[
                \mathcal O_w\text{-span}\,c(\mathcal L_u)=M_{u,w}.
 \tag{GCF.34}
\]
Its \texttt{CmpTied} is now (GCF.33), the common-frame identities,
and these original value/evaluation/character identities, with
\(G_t\)'s original positions.  Thus the third compatibility field
is supplied by a proved map, including its integral span.

For the fourth field the four rings in
\texttt{RingCompatibleAt} are already the canonical
\(WQ(p,n),QQ(p,n),ORec(D,w,n),KfRec(D,w,n)\).
Use the four identity isomorphisms.  They commute with the two
fraction-field embeddings, the maps \(\iota\) and \(j\), and the
specified image of \(w_U\), because these are the maps used in
\(\mathrm{packC}\).  This proves all five commuting equations
for every original \(u,w,v\); no isomorphism with a different
coefficient ring is substituted.  For the fifth field use the actual
\texttt{HidaSetting}, original orbit number, dominant weights and
rank \(k\), the chosen \(I_{u,w},\psi_{u,w}\), and its literal
\(\mathrm{gatOfLevelP}\) equality.  The \texttt{HypAtP} condition
holds because these split local models carry the original self-dual
base lattice by the specified frames, as LI5 and SF6 show.
This is precisely \texttt{GatIs}.

The original recursive witness lists exactly these \(D,L_0,I,\psi,
h_d,B,i_K,\lambda\) and the six-field comparison just constructed.
All five entries in its dependent equality are the definitions in
(GCF.32), so it supplies \texttt{IsRecG}, rather than just equality
of their ranks.  LI6 also supplies, for every original position
\((v,s)\) and every \(x\in\mathcal O_F\),
\[
 x\in\mathfrak p_v
 \ \Longleftrightarrow\
  0<\nu_w\bigl(\tau_{G_t(v,s)}(x)\bigr).             \tag{GCF.35}
\]
It includes \(x=0\) and uses the original normalized valuation and
coefficient embedding.  Finally, the recursive numerical fields are
\[
 r=\begin{cases}r_G(\mathrm{sc})&p<p_0(\mathrm{sc},k),\\1&
                                  p\ge p_0(\mathrm{sc},k),
       \end{cases}\qquad
 e=\left\lceil f k^2\log_p a_{\mathrm{sc}}\right\rceil_+.
 \tag{GCF.36}
\]
These are the two original \texttt{GatNum} equalities, with the
unchanged \(C_e=1\).  This proof constructs the local models jointly
with the comparison.  It does not assert that an arbitrary previously
chosen \(D.\mathrm{gat}\) or \(c\) equals these constructed data.
\hfill\(\square\)
% END GCF12

% BEGIN GCF13
\paragraph{GCF13. The local choices precede every later stage.}
\label{claim:GCF13}
Fix the original \(B,C\), and let \(Z,Z'\) be any two of its
five-part \texttt{Later} records.  All local choices in GCF12 can be
kept when moving from \(Z\) to \(Z'\).  In particular the original
\texttt{IsLocalFamily} for a chosen \(g\) transports to \(Z'\)
with that same \(g\), the same \(B,I,\psi,i_K,\lambda\), and
the transported level-zero forms.

The relevant dependency check is concrete.  The host, coefficient
field, its places and embeddings, ranks, word and class sets, carried
lattices, level-zero groups and characters are fields of the fixed
base and columns.  They do not depend on the later
\(W,V,S,N,G\).  The Frobenius choice \(h_P\), orbit positions,
local weights \(\lambda_{\rm rec}\), and coefficient embedding
\(\tau_{\rm rec}\) read only those fields.  Hence they too are
identical.  The level-zero form transport keeps each underlying
vector space, value function and evaluation map, and merely presents
its level and chart records over the new \(D\).
Every equation of \texttt{IsLocalFamily} therefore has the same two
sides after transport: the standard \(B\), its local frames,
\(i_K\mathrm{embRec}=\lambda\mathrm{embLL}\), all position
identities, the inverse away character, and \(g=\mathrm{gatOf}\).

For the full symmetric columns, the same carrier and form data define
\(\mathrm{colsRe}_{Z'}(C)\).  The endpoint and middle
\texttt{Reads} equations, residue-level identities, and actual value
functions are unchanged, so their \(\mathrm{StdC}\) and every
\(\mathrm{StdLvC}\) remain true.  Taking degree zero commutes
exactly with transport:
\[
  (\mathrm{colsRe}_{Z'}C).\mathrm{forms}(u).\mathrm{to0}
              =\mathrm{forms0Re}_{Z'}(C.\mathrm{forms}(u).\mathrm{to0}).
 \tag{GCF.37}
\]
When only the \texttt{gat} field is replaced by the chosen \(g\),
this is exactly the \(\mathrm{forms0Set}\) used in GCF12.
The value polynomials continue to have the same universal evaluation
formula; VE7's injectivity makes their equality literal, rather than
an unidentified linear equivalence.  Thus the common local choices
need not be repeated after selecting a row stage, division stage,
metric stage, or global group stage.
\hfill\(\square\)
% END GCF13

% BEGIN GCF14
\paragraph{GCF14. One actual \(G\) supplies the full original \(\mathrm{GIs2}\).}
\label{claim:GCF14}
For any original column-choice property \(P\), the literal statement
\[
 \mathrm{StepW}\ P\
  \left(\lambda K,B,f,T,w.\ \exists G:\mathrm{GJStage}(B,w.C),\
                                  \mathrm{GIs2}(K,w,G)\right)
 \tag{GCF.38}
\]
holds.  In particular it holds with \(P=\mathrm{Pinned}\), which
is the original endpoint
\[
 \texttt{exists\_GJStage\_GIs2\_pinned}.
\]
Its witnesses preserve the given coefficient field, embeddings,
column data and ambient period root.

Fix arbitrary objects quantified by \texttt{StepW}, including its
\(K.\mathrm{Over}\) and \texttt{Forced} guards, and a given
\(w:\mathrm{ColW}(B,f,T)\).  The conclusion does not use any
additional clause of \(P\).  The \texttt{Over} data give the
original \texttt{HasLetters}; LS22 and HF19 construct an initial
\(Z_0=(W_0,V_0,S_0,N_0,G_0)\) over these same \(B,w.C\).
TG12 gives its actual \(C_0:\mathrm{ColsTiedSym}\), \(\mathrm{StdC}\),
and all \(\mathrm{StdLvC}\).  Its companion identity is
\[
           L_{b,u}+(\bar\pi_u)L_{b,u}^{\#}=\Lambda_0
\]
by TG12, using the principal ideal generated by the actual
\(\bar\pi_u\).  Its \texttt{Joint0} is part of these columns,
and FC16 supplies the coefficient-field level-zero character with
the original additive-character values.  The local construction in
GCF12 therefore chooses one \(g\) and its exact rank identities
\(h_d\), at scale \(K.s\).

Choose, now and once,
\[
         G=\{G_0.\mathrm{withTrue}
                        \text{ with }\mathrm{gat}=g,
                                      \mathrm{gat\_d}=h_d\}.
 \tag{GCF.39}
\]
The \(\mathrm{withTrue}\) operation keeps its original \(cd\)
and all other stage data, uses the specified true volume, and sets
\(gx=\mathrm{trueGx}\).  GV19 proves \texttt{VolTieG} for it,
with the original Haar normalization and original true-volume
foundation retained there.  Replacing its \texttt{gat} field does
not change that volume identity.

Only after (GCF.39), let \(W,V,S,N\) be arbitrary later stages,
and put \(Z=(W,V,S,N,G)\).  GCF13 transports \(C_0\) and all its
standardness data to this \(Z\), and keeps the chosen local
\(I,\psi,B,i_K,\lambda,g\).  GCF12 then supplies
\((G_t,c,\mathcal G)\), the joint \texttt{IsRecG}, and the exact
prime identity (GCF.35), with the transported forms and the same
level-zero datum.  GCF6 extends the symmetric columns using
\emph{this} \(\mathcal G.\mathrm{herm.g}\) as its prescribed
group-tie field.  Its level and local ties are therefore the two
entries of \(\mathcal G.\mathrm{herm}\), giving the exact equality
\[
 \mathcal G.\mathrm{herm}=F_u.\mathrm{to0},\qquad
 \mathcal G.\mathrm{forms}(u)=C.\mathrm{forms}(u).\mathrm{to0}.
 \tag{GCF.40}
\]
The \texttt{gpl} clause is (GCF.35) on those same ties.  Thus
\(C,\mathrm{StdC},\mathrm{StdLvC},F_u,\mathrm{StdDt},G_t,c,
\mathcal G\), (GCF.40), and \texttt{gpl} are every original
\texttt{ColT} field.

The \texttt{ColW} input supplies \texttt{FieldOK},
\(\mathrm{ccK}=\mathcal O_K\), and the exact twelfth-root law for
\(w.I.J.\delta\).  GCF11 supplies \texttt{OmegaCompatible} for
this \(Z\).  Adding it and the already constructed \texttt{IsRecG}
to \texttt{ColT} gives the entire original
\texttt{ColumnDataSplitPrimes}(K,w,Z).  Finally (GCF.36) gives
\texttt{GatNum}.  We have proved, with \(G\) preceding every
later variable,
\[
 \begin{gathered}
 \forall W,V,S,N,\qquad
       \mathrm{VolTieG}(w.DZ)\ \land\ \mathrm{GatNum}(K.s,w.DZ)\\
       {}\land\ \operatorname{Nonempty}
                  (\mathrm{ColumnDataSplitPrimes}(K,w,Z)).
 \end{gathered}                                                   \tag{GCF.41}
\]
This is exactly \(\mathrm{GIs2}\).  No Bessel nonvanishing,
row-count formula, theta-kernel identity or arbitrary-family
compatibility assumption was inserted into (GCF.38).
\hfill\(\square\)
% END GCF14

% BEGIN GCF15
\paragraph{GCF15. Arbitrary division-lattice changes preserve the comparison map.}
\label{claim:GCF15}
Let \(N_u\) be arbitrary \(\mathcal O_{D.\mathrm{LL}}\)-submodules
of the original division vector spaces and let \(h_u\in\mathbb R\)
be arbitrary.  Let \(D'=\mathrm{setDiv}(D,N,h)\), which changes
only the designated division-lattice and height fields.  Then the
actual recursive comparison satisfies
\[
 \mathrm{cmpTr}_{N,h}(\mathrm{latCmpC}(T,G_t,I,H))
   =\mathrm{latCmpC}(T',G_t',I',H').                  \tag{GCF.42}
\]
All symbols on the right are the original transports, not new
independently chosen forms or comparisons.  No finite-generation,
full-rank or height bound is imposed on the newly inserted \(N_u,h_u\).

The source \texttt{CmpIn} contains the coefficient embedding,
integrality equivalence, common \(B\), its standard-coordinate
identity and the local frame identities.  None of these mentions
the changed fields.  The coefficient polynomial rings, original
form vectors and the local lattice tuples are also unchanged.
Consequently the six \texttt{CmpHyp0P} assertions
\(\mathrm{ex},\mathrm{evInj},\mathrm{ptEx},\mathrm{ptAll},
\mathrm{chi},\mathrm{latM}\) are the same assertions on those
same elements.  Copying them proves the transported hypothesis.
The derived \texttt{CmpHyp} assertions of value existence and
uniqueness, positions, characters, linearity and integral span
transport for the same reason.

For an actual lattice tuple, the map on either side of (GCF.42)
is the coefficient embedding of the unique polynomial satisfying
the VE7 evaluation formula.  The formula and all its arguments are
unchanged, so the polynomials agree.  On a nonlattice tuple both
maps have the source's zero value.  The coefficient embeddings
themselves agree.  These observations prove equality of the full
\texttt{LatCmp} records, rather than only equality of their images
or ranks.

Suppose now \(\mathrm{IsRecG}(\mathrm{sc},D,G_t,c,H_0,F)\).
Take its original recursive witness over \(D_0\).  Replace \(D_0\)
by \(\mathrm{setDiv}(D_0,N,h)\), retaining the same levels,
local identifications, away characters, rank equality, matrix \(B\),
completion embeddings and position identities.  Copy its six-field
comparison hypothesis as above.  Equation (GCF.42) proves the
comparison entry of the required dependent equality; the other four
entries are the original transports.  This gives the literal
\[
 \mathrm{IsRecG}(\mathrm{sc},D',G_t',c',H_0',F').        \tag{GCF.43}
\]
It is the full original
\texttt{concreteSplitPrimeAutomorphicData\_divisionLatticeChange}.
\hfill\(\square\)
% END GCF15

% BEGIN GCF16
\paragraph{GCF16. The complete column compatibility survives that same change.}
\label{claim:GCF16}
The operation of GCF15 also transports the full original
\(\mathrm{GroupCompatibilityAll}\), \(\mathrm{ColT}\), and
\(\mathrm{ColumnDataSplitPrimes}\), keeping their actual
\(G_t,c\) as the specified \(\mathrm{gatTr},\mathrm{cmpTr}\).

The Hermitian spaces, local groups, split-place maps, coefficient
rings, class representatives and form evaluations do not depend on
the changed division submodules.  Therefore the five fields
\(\mathrm{herm},\mathrm{forms},\mathrm{cmp},\mathrm{ring},
\mathrm{gat}\) copy over directly; in the comparison field use
the equality of actual maps in GCF15.  The full Hermitian column
record copies its determinant map and all its laws, every level
and residue map, its volume and relative determinant.  Its forgotten
columns are still the transported original columns.  Every standard
reading, the Hermitian and form equalities, and \texttt{gpl} have
unchanged arguments.  These are all the \texttt{ColT} fields.

For \texttt{ColumnDataSplitPrimes}, (GCF.43) supplies its
\texttt{isRec} field.  The \texttt{om} field also copies: the host,
complex embeddings, \(gx\), the period lattice, its coefficient
data and its one-dimensional different-ideal modules are all
unchanged.  This proves the exact
\texttt{columnLocalDataTransportDivision}, including both extra
fields.  In particular later arithmetic choices of the division
lattice cannot force a fresh choice of a comparison or a fresh
coefficient field.  This statement concerns the named transports;
it is not a construction of every other field of a later
\texttt{Extra} record.
\hfill\(\square\)
% END GCF16

% BEGIN GCF17
\paragraph{GCF17. Equality with the given extra comparison is a joint equality.}
\label{claim:GCF17}
For every original supplied \(q:\mathrm{WPt}(\mathcal S,K)\),
the literal conclusion of \texttt{WPt.extraTied\_col} holds:
\[
 \begin{gathered}
  \mathrm{ColsExtraTiedS}(q.x.\mathrm{col.full},q.O.\mathrm{ex}),\\
  \exists\mathcal G:\mathrm{GroupCompatibilityAll}
                (q.w.D(q.Z),q.O.\mathrm{gatT},q.O.\mathrm{ex.cmp}),\\
  \mathcal G=q.x.\mathrm{col.g},\quad
  \mathcal G.\mathrm{herm}=q.x.\mathrm{col.full.to0},\quad
  \mathcal G.\mathrm{forms}(u)=q.x.\mathrm{col.C.forms}(u).\mathrm{to0}.
 \end{gathered}                                                   \tag{GCF.44}
\]
The equality of \(\mathcal G\) is read after the original dependent
identifications.  The assertion is for the given \(q.O\), including
its supplied comparison and kernel.

The \texttt{Wit} in \(q\) contains the exact equality
\[
 (O.\mathrm{gatT},O.\mathrm{cmp},O.\mathrm{tw},O.\mathrm{ker})
       =(\mathrm{col}.G_t,\mathrm{col}.c,w.\mathrm{tw},\mathrm{ker}_0)
 \tag{GCF.45}
\]
in the dependent space in which the kernel is indexed by the first
three entries.  Substitute this equality once, simultaneously.  Use
\(\mathrm{col}.g\) as \(\mathcal G\); its existing Hermitian and
form equalities are now precisely the last two equations in
(GCF.44).  The \texttt{ColsExtraTiedS} level and error-datum fields
are \(\mathrm{Wit.lev}\) and \(\mathrm{Wit.errX}\), while its
prime compatibility is \(\mathrm{col.gpl}\).  Its two kernel
embedding clauses are the actual \(\mathrm{KerT.iw}\) and
\(\mathrm{KerT.iw\_pos}\), transported through (GCF.45).
Thus every field is on the same given kernel and comparison.

The resulting family implication has an exact boundary.  Suppose a
given family is eventually realized by such \(q\)'s with the
original joint equality of its \((D,O)\), for every family index
and every later admissible height witness.  Then (GCF.44) gives
\(\mathrm{GroupCompatibilityAll}(\mathrm{Fm}.D,
\mathrm{Fm}.\mathrm{gatT},\mathrm{Fm}.\mathrm{ex.cmp})\) at
those same points.  Taking the maximum of the realization threshold
and \(\mathrm{Fm}.x_1+1\) gives exactly
\(\mathrm{GroupCompatibilityFamily}\).  Likewise the original
\texttt{TSelF}, if supplied, provides its actual \texttt{RecT}
at every point, whose \(\mathrm{TiedDataS.g}\) is already the
same compatibility record; the same maximum gives the original
\texttt{gTiedF\_of\_sel}.

GCF14 constructs the complete column input used by these witnesses,
and GCF16 preserves it through division-lattice choices.  Neither
of those conclusions alone supplies the other \texttt{Wit} fields:
the selected rows, kernel/theta and divisor ties, objects, gauge,
and the joint equality (GCF.45) must still be realized in the common
construction.  Their existence is an explicit remaining producer,
not a consequence of choosing separate compatible records at the
individual places.
\hfill\(\square\)
% END GCF17

% BEGIN GCF18
\paragraph{GCF18. The selected design class already pays for its top field.}
\label{claim:GCF18}
For the original class
\[
 D_3(C_0,\mathcal A)=\mathrm{Des}(\mathrm{WithChN}(\mathcal A),C_0),
\]
membership supplies \(\exists C_H\,\mathrm{ChosenE}(\mathrm{Fm},C_0,C_H)\).
That record contains the actual selected \(\mathrm{ChosenTT}\)
and one positive constant \(C_b\) for \(\mathrm{ChosenBad}\).
Hence it supplies \(\mathrm{TopsF}\) on the unchanged selected
subset.  GLG16 proves its full original \(\mathrm{Adm2F}\), with
the unchanged explicit threshold
\[
 x_0=\max\{x_a,x_b,\mathrm{Fm}.x_1+1,
                              (16C_b+2)^{3/\epsilon}\}.
 \tag{GCF.46}
\]
No new selected subset or independent probability experiment occurs.
In particular \(\mathrm{Adm2F}\) is no longer a separate producer
once membership in this selected design class is given.

If the original selector \(\mathrm{TSel}(\mathcal A)\) is
supplied, GCF17 gives the first \texttt{Closed} input required by
LGS13 on \(D_3(3,\mathcal A)\).  The preceding projection gives
the second \texttt{Closed} input for \(\mathrm{Adm2F}\).
These are class-level statements: they hold for every family in the
class and are available before the later family and height choices.
Take the pointwise maximum of their original coefficient bounds and
the finite bounds in LGS13.  Apply that claim on the same family.
This yields exactly the original
\[
 \mathrm{Closed}\bigl(D_3(3,\mathcal A),\mathrm{LinkG}\bigr),
 \tag{GCF.47}
\]
with all five \(\mathrm{LinkG}\) fields, its error
\(b_o=b^{x^{-\epsilon/2}}\), and its unchanged
\(C_G=4C_{\rm count}+1\).  The original
\texttt{linkG\_frame} follows with its literal subsequent
\(D_3,\mathrm{TuSize},\mathrm{DGamSize}\) guards; it introduces
no stronger coefficient quantifier.

The antecedent \(\mathrm{TSel}(\mathcal A)\), or an exact eventual
realization sufficient for GCF17, is retained in (GCF.47).
The completed supplier in this chapter is the full earlier
\(\mathrm{GIs2}\), its determinant and period inputs, and its
literal comparison transport.  No whole \texttt{GLOBAL\_LOCAL\_GAIN}
closure, no other global quantitative leaf, and no complete RH or GRH
proof is asserted here.
\hfill\(\square\)
% END GCF18
